% ============================================================
%  Divisibilidade — Apostila Premium | PratiqueJa
%  Engine: XeLaTeX | Font: Inter
% ============================================================
\documentclass[12pt]{article}

\usepackage[brazil]{babel}
\usepackage{fontspec}
\setmainfont[
  BoldFont=Inter-Bold, ItalicFont=Inter-Italic,
  BoldItalicFont=Inter-BoldItalic,
]{Inter-Regular}
\usepackage{microtype}

\usepackage[a4paper, top=1.5cm, bottom=2.8cm,
            left=1.5cm, right=1.5cm]{geometry}
\setlength{\footskip}{1.3cm}

% ── Paleta de cores ───────────────────────────────────────────
\usepackage[dvipsnames,svgnames,table]{xcolor}

\definecolor{darkblue}{HTML}{1E3A8A}     % azul principal
\definecolor{accentblue}{HTML}{38BDF8}   % azul acento
\definecolor{graycolor}{HTML}{64748B}    % cinza neutro
\definecolor{vegreen}{HTML}{22C55E}      % verde correto
\definecolor{vered}{HTML}{EF4444}        % vermelho incorreto
\definecolor{badgepurple}{HTML}{7C3AED}  % roxo — definição
\definecolor{badgegreen}{HTML}{059669}   % verde — avançado

\definecolor{rulebg}{HTML}{EEF2FF}       % fundo regra (azul muito suave)
\definecolor{exbg}{HTML}{F8FAFC}         % fundo exemplos
\definecolor{obsbg}{HTML}{F1F5F9}        % fundo observação (cinza suave)
\definecolor{testebg}{HTML}{FFFBEB}      % fundo teste rápido (âmbar suave)
\definecolor{exborder}{HTML}{E2E8F0}     % borda neutra
\definecolor{testebd}{HTML}{FDE68A}      % borda teste (âmbar)
\definecolor{sepcolor}{HTML}{E2E8F0}     % linhas divisórias
\definecolor{bodytext}{HTML}{374151}     % texto principal
\definecolor{titletext}{HTML}{0F172A}    % títulos
\definecolor{mutedtext}{HTML}{64748B}    % texto secundário

% ── Pacotes ───────────────────────────────────────────────────
\usepackage{amsmath,amssymb}
\usepackage{multicol}
\usepackage{setspace}
\usepackage{parskip}
\usepackage{array}

\usepackage[most]{tcolorbox}
\tcbuselibrary{skins, breakable}

% ── Rodapé ───────────────────────────────────────────────────
\usepackage{fancyhdr}
\pagestyle{fancy}
\fancyhf{}
\renewcommand{\headrulewidth}{0pt}
\renewcommand{\footrulewidth}{0.5pt}
\renewcommand{\footrule}{%
  {\color{sepcolor}%
   \vskip-\footruleskip\vskip-\footrulewidth%
   \hrule width\headwidth height\footrulewidth%
   \vskip\footruleskip}}
\fancyfoot[L]{\small\color{mutedtext}%
  \textbf{Pratique}\textcolor{darkblue}{\textbf{Já}}}
\fancyfoot[C]{\small\color{mutedtext}\itshape Divisibilidade}
\fancyfoot[R]{\small\color{mutedtext}\thepage}

% ============================================================
%  ATALHOS TIPOGRÁFICOS
% ============================================================
\newcommand{\hl}[1]{\textcolor{darkblue}{\ifmmode\mathbf{#1}\else\textbf{#1}\fi}}
\newcommand{\ok}{{\bfseries\textcolor{vegreen}{✓}}}
\newcommand{\nok}{{\bfseries\textcolor{vered}{✗}}}

% ============================================================
%  COMPONENTES VISUAIS
% ============================================================

% ── Cabeçalho da apostila ─────────────────────────────────────
%    Linha azul acento no topo · título · marca · subtítulo · separador
\newcommand{\heroheader}[3]{%
  \noindent{\color{accentblue}\rule{\linewidth}{2.5pt}}\par
  \vspace{7pt}%
  \noindent
  {\color{titletext}\fontsize{22}{27}\selectfont\bfseries #1}%
  \hfill
  {\color{mutedtext}\small
    \textbf{Pratique}\textcolor{darkblue}{\textbf{Já}}%
    \ \textcolor{exborder}{|}\ %
    \itshape #3}\par
  \vspace{2pt}%
  {\color{mutedtext}\small\itshape #2}\par
  \vspace{6pt}%
  \noindent{\color{sepcolor}\rule{\linewidth}{0.5pt}}%
  \vspace{8pt}%
}

% ── Badge inline ──────────────────────────────────────────────
\newcommand{\badge}[2]{%
  \tcbox[
    on line, enhanced, arc=5pt, boxrule=0pt,
    colback=#2, colframe=#2,
    left=6pt, right=6pt, top=3pt, bottom=3pt,
    nobeforeafter,
  ]{\footnotesize\bfseries\color{white}#1}%
}

% ── Cabeçalho de seção ────────────────────────────────────────
%    Linha divisória suave · badge + título · subtítulo
\newcommand{\TW}[4]{%
  \vspace{7pt}%
  \noindent{\color{graycolor!35}\rule{\linewidth}{0.5pt}}\par\nopagebreak
  \vspace{4pt}%
  \noindent\badge{#1}{#2}\hspace{6pt}%
  {\color{titletext}\bfseries\normalsize #3}%
  {\color{mutedtext}\ $\cdot$\ \small\itshape #4}\par\nopagebreak
  \vspace{5pt}%
}

% ── Regra — definição/regra principal ────────────────────────
%    Fundo azul suave · borda esquerda escura · texto em bodytext
\newcommand{\regra}[1]{%
  \begin{tcolorbox}[
    enhanced, arc=3pt, boxrule=0pt,
    colback=rulebg, colframe=rulebg,
    left=10pt, right=8pt, top=6pt, bottom=6pt,
    borderline west={3pt}{0pt}{darkblue},
  ]
    \color{bodytext}\small #1
  \end{tcolorbox}%
  \vspace{2pt}%
}

% ── Exemplos agrupados — caixa única compacta ─────────────────
\newcommand{\exemp}[1]{%
  \vspace{2pt}%
  \begin{tcolorbox}[
    arc=5pt, boxrule=0.5pt,
    colback=exbg, colframe=exborder,
    left=9pt, right=9pt, top=5pt, bottom=5pt,
  ]
    \setlength{\parskip}{1.5pt}%
    \small\color{bodytext}#1
  \end{tcolorbox}%
  \vspace{1pt}%
}

% ── Fórmula destacada ─────────────────────────────────────────
%    Linha acento azul no topo · fundo suave · borda escura
\newcommand{\formula}[1]{%
  \vspace{3pt}%
  \begin{tcolorbox}[
    enhanced, arc=5pt, boxrule=0.7pt,
    colback=rulebg, colframe=darkblue,
    left=10pt, right=10pt, top=8pt, bottom=8pt,
    borderline north={2pt}{0pt}{accentblue},
  ]
    \centering\color{darkblue}\normalsize\bfseries #1
  \end{tcolorbox}%
  \vspace{3pt}%
}

% ── Observação — nota discreta com borda cinza ────────────────
\newcommand{\obs}[1]{%
  \vspace{2pt}%
  \begin{tcolorbox}[
    enhanced, arc=3pt, boxrule=0pt,
    colback=obsbg, colframe=obsbg,
    left=8pt, right=8pt, top=4pt, bottom=4pt,
    borderline west={2pt}{0pt}{graycolor!60},
  ]
    {\color{graycolor}\fontsize{7}{8.5}\selectfont\bfseries obs.\enspace}%
    \small\color{bodytext}#1
  \end{tcolorbox}%
  \vspace{1pt}%
}

% ── Teste rápido — exercício compacto em âmbar suave ──────────
\newcommand{\teste}[1]{%
  \vspace{2pt}%
  \begin{tcolorbox}[
    arc=4pt, boxrule=0.5pt,
    colback=testebg, colframe=testebd,
    left=8pt, right=8pt, top=3pt, bottom=3pt,
  ]
    {\color{graycolor}\fontsize{7}{8.5}\selectfont\bfseries
     teste rápido\enspace}%
    \small\color{bodytext}#1
  \end{tcolorbox}%
  \vspace{1pt}%
}

% ── Rótulo interno (h4) ───────────────────────────────────────
\newcommand{\SH}[1]{%
  \vspace{4pt}%
  \noindent{\color{mutedtext}\fontsize{7.5}{9}\selectfont
            \bfseries\MakeUppercase{#1}}\par
  \vspace{1pt}%
}

% ── Chips de divisor ──────────────────────────────────────────
\newcommand{\divisorchip}[1]{%
  \tcbox[
    on line, enhanced, arc=5pt, boxrule=0.5pt,
    colback=rulebg, colframe=exborder,
    left=4pt, right=4pt, top=2pt, bottom=2pt,
    nobeforeafter,
  ]{\small\bfseries\color{darkblue}#1}\hspace{2pt}%
}

% ============================================================
\begin{document}
\pagenumbering{arabic}
\setstretch{1.2}
\color{bodytext}

\heroheader{Divisibilidade}%
           {Regras, Divisores e Aplicações}%
           {Matemática · Intermediário}

\setlength{\columnsep}{18pt}
\setlength{\columnseprule}{0.4pt}
\def\columnseprulecolor{\color{accentblue!30}}

\begin{multicols}{2}

%─────────────────────────────────────────────────────────────
\TW{Def}{badgepurple}{Definição}{Divisibilidade}

\regra{Dizemos que $a$ é \textbf{divisível} por $b$ quando existe
$q \in \mathbb{Z}$ tal que $a = b \cdot q$, ou seja, o \textbf{resto
da divisão é zero}. Chamamos $b$ de \textbf{divisor} de $a$ e $a$
de \textbf{múltiplo} de $b$.}

%─────────────────────────────────────────────────────────────
\TW{2}{darkblue}{Divisibilidade por 2}{Números pares}

\regra{O \textbf{último dígito} deve ser \textbf{0, 2, 4, 6 ou 8}.}

\exemp{%
  13\hl{8} → último dígito 8 → divisível \ok\\
  47\hl{5} → último dígito 5 → não divisível \nok%
}
\teste{a) 734 \quad b) 519 \quad c) 2.460}

%─────────────────────────────────────────────────────────────
\TW{3}{darkblue}{Divisibilidade por 3}{Soma dos algarismos}

\regra{A \textbf{soma dos algarismos} deve ser múltipla de \textbf{3}.}

\exemp{%
  $132 \Rightarrow 1{+}3{+}2=\hl{6}$ → múltiplo de 3 \ok\\
  $531 \Rightarrow 5{+}3{+}1=\hl{9}$ → múltiplo de 3 \ok\\
  $124 \Rightarrow 1{+}2{+}4=\hl{7}$ → não múltiplo de 3 \nok%
}
\teste{a) 513 \quad b) 724 \quad c) 3.006}

%─────────────────────────────────────────────────────────────
\TW{4}{darkblue}{Divisibilidade por 4}{Dois últimos dígitos}

\regra{Os \textbf{2 últimos dígitos} devem formar número divisível
por 4. Os demais algarismos não influenciam.}

\exemp{%
  $5\hl{24}$ → 24 é múltiplo de 4 \ok\\
  $13\hl{36}$ → 36 é múltiplo de 4 \ok\\
  $5\hl{14}$ → 14 não é múltiplo de 4 \nok%
}
\teste{a) 748 \quad b) 326 \quad c) 1.512}

%─────────────────────────────────────────────────────────────
\TW{5}{darkblue}{Divisibilidade por 5}{Último dígito 0 ou 5}

\regra{O \textbf{último dígito} deve ser \textbf{0} ou \textbf{5}.}

\exemp{%
  $32\hl{5}$ → termina em 5 \ok\quad
  $14\hl{0}$ → termina em 0 \ok\\
  $78\hl{3}$ → termina em 3 → não divisível \nok%
}
\teste{a) 345 \quad b) 472 \quad c) 1.000}

%─────────────────────────────────────────────────────────────
\TW{6}{darkblue}{Divisibilidade por 6}{Divisível por 2 e por 3}

\regra{Deve ser \textbf{par} e a soma dos algarismos deve ser
\textbf{múltipla de 3} — ambas as condições simultaneamente.}

\exemp{%
  $654$: par e $6{+}5{+}4=\hl{15}$ ($\div3$) \ok\\
  $312$: par e $3{+}1{+}2=\hl{6}$ ($\div3$) \ok\\
  $125$: ímpar → não divisível por 6 \nok%
}
\obs{Basta uma condição falhar para o número não ser divisível por 6.}
\teste{a) 432 \quad b) 514 \quad c) 3.612}

%─────────────────────────────────────────────────────────────
\TW{7}{darkblue}{Divisibilidade por 7}{Regra do duplo}

\regra{Multiplique o \textbf{último dígito por 2} e subtraia do
número sem esse dígito. Repita até concluir.}

\exemp{%
  $16\hl{1} \Rightarrow 1{\cdot}2=2 \Rightarrow 16{-}2=\hl{14}$ \ok\\
  $109\hl{9} \Rightarrow 9{\cdot}2{=}18 \Rightarrow 109{-}18{=}91
  \Rightarrow 9{-}2{=}\hl{7}$ \ok%
}
\obs{Para números pequenos, a divisão direta costuma ser mais prática.}
\teste{a) 203 \quad b) 140 \quad c) 511}

%─────────────────────────────────────────────────────────────
\TW{8}{darkblue}{Divisibilidade por 8}{Três últimos dígitos}

\regra{Os \textbf{3 últimos dígitos} devem formar número divisível
por 8. Os demais não influenciam.}

\exemp{%
  $7\hl{128}$ → 128 é múltiplo de 8 \ok\\
  $5\hl{256}$ → 256 é múltiplo de 8 \ok\\
  $3\hl{100}$ → 100 não é múltiplo de 8 \nok%
}
\teste{a) 1.024 \quad b) 516 \quad c) 3.112}

%─────────────────────────────────────────────────────────────
\TW{9}{darkblue}{Divisibilidade por 9}{Soma dos algarismos}

\regra{A \textbf{soma dos algarismos} deve ser múltipla de \textbf{9}.}

\exemp{%
  $432 \Rightarrow 4{+}3{+}2=\hl{9}$ \ok\quad
  $124 \Rightarrow 1{+}2{+}4=\hl{7}$ \nok\\
  $6561 \Rightarrow 6{+}5{+}6{+}1=\hl{18}$ → múltiplo de 9 \ok%
}
\obs{Todo número divisível por 9 também é divisível por 3.}
\teste{a) 729 \quad b) 324 \quad c) 1.035}

%─────────────────────────────────────────────────────────────
\TW{10}{darkblue}{Divisibilidade por 10}{Último dígito 0}

\regra{O \textbf{último dígito deve ser 0}. Como $10=2{\times}5$,
equivale a ser divisível por 2 e por 5 simultaneamente.}

\exemp{%
  $15\hl{0}$ → termina em 0 \ok\quad
  $487\hl{0}$ → termina em 0 \ok\\
  $32\hl{5}$ → termina em 5 → divisível por 5, mas não por 10 \nok%
}
\teste{a) 1.200 \quad b) 375 \quad c) 4.800}

%─────────────────────────────────────────────────────────────
\TW{11}{badgegreen}{Divisibilidade por 11}{Soma alternada dos algarismos}

\regra{A \textbf{soma alternada} dos algarismos (alternando $+$ e
$-$ da esquerda para a direita) deve ser múltipla de 11 ou zero.}

\exemp{%
  $253 \Rightarrow {+}2{-}5{+}3=\hl{0}$ → múltiplo de 11 \ok\\
  $1452 \Rightarrow {+}1{-}4{+}5{-}2=\hl{0}$ \ok\\
  $123 \Rightarrow {+}1{-}2{+}3=\hl{2}$ → não divisível \nok%
}
\teste{a) 198 \quad b) 352 \quad c) 1.001}

%─────────────────────────────────────────────────────────────
\TW{Div}{badgegreen}{Divisores de um Número}{Decomposição em fatores primos}

Para encontrar os divisores: decomponha em fatores primos e use a
fórmula abaixo para \textbf{contar} sem listar todos.

\SH{Decomposição — 60}

\vspace{2pt}
\noindent
$\begin{array}[t]{r@{\ }|@{\ }l}
60 & 2 \\ 30 & 2 \\ 15 & 3 \\ 5 & 5 \\ \hline 1 &
\end{array}$
\quad$\Rightarrow\ 60 = 2^2 \cdot 3^1 \cdot 5^1$

\SH{Divisores de 60}

\vspace{2pt}\noindent
\divisorchip{1}\divisorchip{2}\divisorchip{3}\divisorchip{4}%
\divisorchip{5}\divisorchip{6}\divisorchip{10}\divisorchip{12}%
\divisorchip{15}\divisorchip{20}\divisorchip{30}\divisorchip{60}

\SH{Fórmula — quantidade de divisores}

Se $n = p_1^{a_1} \cdot p_2^{a_2} \cdots$, então:

\formula{$d(n) = (a_1+1)\cdot(a_2+1)\cdots$}

\exemp{%
  $d(60)=(2{+}1)(1{+}1)(1{+}1)=3{\cdot}2{\cdot}2=\hl{12}$ divisores\\
  $d(36)=(2{+}1)(2{+}1)=3{\cdot}3=\hl{9}$ divisores%
}
\teste{Quantos divisores: a)~$12$ \quad b)~$48$ \quad c)~$100$}

\end{multicols}

\end{document}
